\section{Datasets and experimental design}
\subsection{Two operational adapters}
\paragraph{Compute cluster.}
We use all 500 task-event shards of Google ClusterData2011-2
\citep{google2011}. The trace covers approximately 29 days in one cluster.
Positive trace timestamps use the publisher's 600-second physical origin;
time-zero rows describe an initial state with unknown preceding age. The
adapter replays observed events in timestamp and source order. Tasks are
inactive, pending, or running. The prediction target is the first change of
the current pending/running state, at 1, 5, 15, and 60 minutes. It is not
necessarily task completion: a running task may fail, be evicted, or be killed.
Current job membership and known running-machine membership define two
relation types. A pending or unknown-machine task has no shared machine factor.
Missing requests retain known preceding values, with missingness indicators;
publisher imputation and unknown initial age are retained as cutoff features.
Requests are not measurements of resource utilization.

\paragraph{Loan application workflow.}
BPI Challenge 2017 records applications initiated during 2016 with later
follow-up \citep{vandongen2017}. We predict the first occurrence of
\texttt{A\_Pending}, \texttt{A\_Denied}, or \texttt{A\_Cancelled} after creation,
at 1, 3, 7, and 14 days. ``Resolution'' denotes this declared milestone; it
does not mean repayment or necessarily the last logged event. Cases with a
qualifying prior resolution are no longer members of the active cohort.
The relation is the last observed workflow handler, updated from supported
workflow lifecycle events. This is historical association, not a claim that
all associated cases exclusively occupy the handler's capacity. Retrospective
trace attributes, eventual offer outcomes, and amounts without an availability
audit are excluded. Activities and lifecycle categories unavailable in the
training vocabulary map to an unknown token.

\subsection{Chronological partitions}
All features are built from events at or before the forecast origin. Labels
are attached afterward. Entity identifiers are storage keys, not learned
embeddings. Ongoing entities carry their observed prefix across partitions;
their future outcomes are never used as prefix features. The longest horizon
must lie within the relevant label window.

For Borg, individual-model training ends at elapsed day 18, calibration uses
days 18--20, and development targets use days 20--22. Dependence fitting uses
forward predictions in days 8--10, 12--14, and 16--18; the corresponding
individual models train before days 6, 10, and 14, with the following two days
reserved for calibration. Final forecast origins are every six hours from
day 22 through day 28.75: 28 complete current cohorts.

For BPI, individual-model training ends before September 2016, calibration
uses September, and development targets use October. Forward predictions for
dependence fitting use June, July, and August, each with an earlier training
and calibration period. Final forecast origins are noon UTC each day from
November 1 through December 31: 61 complete current cohorts. Follow-up through
January 15 supports all 14-day labels. The archive concerns 2016 applications;
this experiment does not require predicting unobserved 2017 arrivals.

The final evaluation contains 3,760,001 task--origin queries in Borg and
114,051 case--origin queries in BPI. Borg cohorts range from 105,077 to 171,603
active tasks; BPI cohorts range from 1,484 to 2,189 active cases. These counts
include repeated entities at different origins and are not independent sample
sizes. The final Borg prefix contains 144,648,283 task-event rows.

Model selection was exploratory and development-informed. Before opening final
labels, we froze model choices, checkpoints, code hashes, origins, metrics,
random seeds, and uncertainty calculations in an immutable manifest. No model
is retrained or recalibrated on final outcomes. This is a development-informed
holdout freeze, not prospective preregistration of the initial research idea.

\subsection{Individual models and common-CDF controls}
Histogram gradient boosting (HGB) uses personal, global, and relational feature
arms. The personal arm describes observed state, age, event history, and
available requests or offer history. The global arm adds calendar features
and system-wide activity; the relational arm additionally summarizes current
job/machine populations or last-handler case history. These arms are nested.
Borg has 32 base features and an extension with 144 recent-event counts;
BPI has 103 features. Four monotone calibration variants are compared: sorted
raw horizon probabilities, independently Platt-calibrated and sorted
probabilities, and two shared positive-slope Platt maps with population or
equal-origin fitting weights. Temporal encoders receive the same extended
tabular features and the last 16 observed events per scope: task, job, and
machine for Borg; case and handler for BPI. Entity identifiers are excluded
from learned embeddings.
The temporal Mixer and GRU use a five-category first-exit head; the MLP uses
the tabular features alone. We compare raw and temperature-calibrated
three-seed mean predictions. Selection minimizes equal-origin integrated
Brier score, separately by current state, without choosing the best seed.
Here integrated Brier means the arithmetic mean of binary Brier scores at
the four specified horizons, with equal forecast-origin weights; it is a
discrete-horizon summary rather than a continuous-time integral.

All dependence controls share this selected CDF. Each family is fitted with
seeds 0, 1, and 2, and evaluated with 512 draws per origin. Seed 0 is also
evaluated with 1,024 draws as a Monte Carlo sensitivity. Both families and all
controls remain in the final tables even when development selection favors a
single configuration. Hyperparameters and complete development comparisons
are reported in the appendix and machine-readable manifests.

\subsection{Historical aggregate-error comparator}
For each development origin and eligible group, let
$\mu_{Ak}=\sum_{i\in A}F_{ik}$ and
$s_{Ak}=\sqrt{1+\sum_{i\in A}F_{ik}(1-F_{ik})}$.
We store the four-dimensional residual
$e_A=(C_A-\mu_A)/s_A$. A draw selects one historical origin for all groups,
then a residual vector from a group with the closest
$\lfloor\log_2 |A|\rfloor$ size bin within the same projection. Population
states retain their own residual vectors. Predicted counts
$\mu_A+s_Ae_A$ are clipped to $[0,|A|]$, rounded, and made nondecreasing across
horizons. Development evaluation leaves out the scored origin; final
evaluation uses all development residuals with no further fitting.
This comparator can change group means and uses development labels as fitting
data. Its separate count projections need not be realizable as a single joint
entity scenario. Accordingly, it is an aggregate-forecast benchmark, not a
matched-marginal dependence ablation.
